\section{模型假设}

结合题面已给条件，仅对题面未完全规定、而建模必须补足的现实简化作如下假设：

\begin{enumerate}
\item[(1)] \textbf{一维轴对称}。因 $L=25\,\mathrm{cm}\gg R_0=2\,\mathrm{cm}$，端面散失相对侧面
      可忽略，温度与含水率只沿半径 $r$ 变化，$\partial_\theta(\cdot)=\partial_z(\cdot)=0$。
      这使三维问题降为一维，保留了径向输运这一主导机制。
\item[(2)] \textbf{连续介质与局部平衡}。药材视为均质各向同性多孔介质，同一微元内固相、
      液态水与热量处于局部热力学平衡，可用单一温度 $T$ 与干基含水率 $C$ 描述其状态。
\item[(3)] \textbf{Fick 型扩散与 Fourier 型导热}。内部水分迁移服从以 $C$ 为势的 Fick 第二定律，
      热量传递服从 Fourier 定律；扩散系数 $D$ 与导热系数 $k$ 按各附录设定随 $C,T$ 变化。
\item[(4)] \textbf{对流型表面边界}。表面失水与散热速率正比于表面与热风之间的含水率差、温差，
      即 Robin 边界，传质、传热系数 $h_m,h$ 取为常数；这刻画了外部对流阻力的作用\upcite{giner2009}。
\item[(5)] \textbf{潜热不单列}。题面未给蒸发潜热 $L$，基准模型不显式扣除表面蒸发吸热，
      将其影响留待灵敏度分析定量评估（见第\ref{sec:sens}节），以避免引入未给定参数。
\item[(6)] \textbf{纯径向均匀收缩}（仅问题四）。附件2 的体积收缩按半径各向同性缩小处理，
      干物质总量守恒，收缩速率由附件2 数据经单调插值确定，不额外假设收缩动力学。
\end{enumerate}

\section{符号说明}

下表收录跨章节反复使用的核心符号，一次性下标与临时派生量在首次出现处就地说明。

\begin{center}
\small
\begin{longtable}{cll}
\toprule
符号 & 含义 & 单位 \\
\midrule
\endfirsthead
\toprule
符号 & 含义 & 单位 \\
\midrule
\endhead
\bottomrule
\endfoot
$r,\ t$ & 径向坐标、时间 & $\mathrm{m},\ \mathrm{s}$ \\
$R_0,\ R(t)$ & 初始半径、$t$ 时刻半径（问题四移动边界） & $\mathrm{m}$ \\
$T,\ T_{\mathrm{air}}$ & 药材温度、热风温度 & $^{\circ}\mathrm{C}$ \\
$C,\ C_{\mathrm{air}}$ & 干基含水率、空气等效含水率 & 1 \\
$C_0,\ C_{\mathrm{th}}$ & 初始含水率、安全含水率阈值 & 1 \\
$C_s,\ T_s$ & 表面含水率、表面温度 & 1,\ $^{\circ}\mathrm{C}$ \\
$\rho,\ c_p$ & 湿基密度、比热容 & $\mathrm{kg\,m^{-3}},\ \mathrm{J\,kg^{-1}K^{-1}}$ \\
$k$ & 导热系数 & $\mathrm{W\,m^{-1}K^{-1}}$ \\
$D$ & 水分扩散系数 & $\mathrm{m^2\,s^{-1}}$ \\
$h,\ h_m$ & 对流换热系数、对流传质系数 & $\mathrm{W\,m^{-2}K^{-1}},\ \mathrm{m\,s^{-1}}$ \\
$T_K$ & 绝对温度 $T+273.15$ & $\mathrm{K}$ \\
$\alpha$ & 热扩散率 $k/(\rho c_p)$ & $\mathrm{m^2\,s^{-1}}$ \\
$\mathrm{Bi}_m$ & 传质毕奥数 $h_m R_0/D$ & 1 \\
$\Phi(C)$ & 扩散系数的 Kirchhoff 积分势 & $\mathrm{m^2\,s^{-1}}$ \\
$\xi$ & 固定域无量纲坐标 $r/R(t)$ & 1 \\
$\sigma$ & 干基质量不变量 $\rho_s R^2$ & $\mathrm{kg\,m^{-1}}$ \\
$w_j$ & 控制体柱坐标测度权 & 1 \\
$t^{*}$ & 内部各点首次全部达标的时刻 & $\mathrm{h}$ \\
\end{longtable}
\end{center}
